\begin{proof}
Write
\[
                  \ell_{ij}(r,s)=c+A(rv_i)+B(sw_j).
\]
For each pair $(i,j)$, the convex function $u_{ij}-\ell_{ij}$ vanishes on
a nonempty open subset of the rectangle.  The zero linear form is therefore
a subgradient there, and convexity gives
\begin{equation}
                         \ell_{ij}\leq u_{ij}
                         \quad\text{on }[0,R_0]^2.
 \label{eq:p071-cross-collar-affine-minorant}
\end{equation}

Choose a finite affine maximum representation
\begin{equation}
                  u_{ij}(r,s)=\max_{1\leq k\leq q_{ij}}
                                  L_{ij,k}(r,s)
 \label{eq:p071-cross-collar-polyhedral-maximum}
\end{equation}
on the rectangle.  Every $L_{ij,k}$ is an affine minorant of $u_{ij}$
there.  In particular, on the cross-collar of
\eqref{eq:p071-cross-collar-affine-data},
\begin{equation}
                         L_{ij,k}\leq\ell_{ij}.
 \label{eq:p071-cross-collar-piece-gap}
\end{equation}

We use an exact regularized maximum.  Fix a nonnegative even
$C^\infty$ function $\theta$ on $\mathbf R$, supported in $[-1,1]$ and
having integral one.  For $q\geq1$ and $\varepsilon>0$, let
\[
 \rho_{\varepsilon,q}(t)
   =\prod_{k=1}^q\varepsilon^{-1}\theta(t_k/\varepsilon),
 \qquad
 M_{\varepsilon,q}(x)
   =\int_{\mathbf R^q}\max_k(x_k+t_k)
      \rho_{\varepsilon,q}(t)\,\mathrm dt.
\]
Then $M_{\varepsilon,q}$ is smooth, convex, and nondecreasing in every
coordinate, and
\begin{align}
 \max_kx_k&\leq M_{\varepsilon,q}(x)
       \leq\max_kx_k+\varepsilon,
 \label{eq:p071-cross-collar-regularized-max-bound}\\
 M_{\varepsilon,q}(x)&=x_a
 \quad\text{if}\quad x_a>x_k+2\varepsilon\quad(k\ne a).
 \label{eq:p071-cross-collar-regularized-max-exactness}
\end{align}
Indeed, convolution preserves convexity and coordinatewise monotonicity;
the lower bound is Jensen's inequality, the upper bound follows from the
support of the mollifier, and the last assertion follows because the
$a$-th entry remains the unique maximum throughout that support.

Put
\[
 C_n=4^{-n},\qquad \varepsilon_n=\frac14C_n,
 \qquad e_n=C_n+\varepsilon_n,
 \qquad d_n=\sum_{k>n}e_k.
\]
Thus
\begin{equation}
                         d_n=e_{n+1}+d_{n+1}.
 \label{eq:p071-cross-collar-tail-identity}
\end{equation}
On branch rectangle $(i,j)$ define
\begin{equation}
 g_{n,ij}=M_{\varepsilon_n,q_{ij}+1}
   \bigl(\ell_{ij}+C_n,L_{ij,1},\ldots,L_{ij,q_{ij}}\bigr).
 \label{eq:p071-cross-collar-smoothed-maximum}
\end{equation}
Equations \eqref{eq:p071-cross-collar-affine-minorant},
\eqref{eq:p071-cross-collar-polyhedral-maximum}, and
\eqref{eq:p071-cross-collar-regularized-max-bound} give
\begin{equation}
                   u_{ij}\leq g_{n,ij}\leq u_{ij}+e_n
                   \quad\text{on }[0,R_0]^2.
 \label{eq:p071-cross-collar-uniform-sandwich}
\end{equation}

On the cross-collar, \eqref{eq:p071-cross-collar-piece-gap} and
$C_n>2\varepsilon_n$ show that the first entry of
\eqref{eq:p071-cross-collar-smoothed-maximum} dominates every other entry
by more than $2\varepsilon_n$.  Hence
\eqref{eq:p071-cross-collar-regularized-max-exactness} gives
\begin{equation}
                         g_{n,ij}=\ell_{ij}+C_n
                         \quad\text{if }\min\{r,s\}<\delta.
 \label{eq:p071-cross-collar-exact-collar}
\end{equation}
The nine functions therefore glue smoothly across both node axes.  Every
branch restriction is jointly convex because it is a convex nondecreasing
function of affine functions.  Moreover, on the first node axis,
\[
 \sum_{i=1}^3\partial_rg_{n,ij}(0,s)
   =\sum_{i=1}^3 A(v_i)=A(v_1+v_2+v_3)=0,
\]
and the analogous identity with $B$ proves the second balance.  Thus the
glued function $g_n$ belongs to $\mathcal G_0(V)$.

Finally set
\[
                              G_n=g_n+d_n.
\]
By \eqref{eq:p071-cross-collar-uniform-sandwich} and
\eqref{eq:p071-cross-collar-tail-identity},
\[
 G_{n+1}\leq u+e_{n+1}+d_{n+1}
            =u+d_n\leq g_n+d_n=G_n.
\]
The same sandwich gives
\[
                 0\leq G_n-u\leq e_n+d_n\longrightarrow0
\]
uniformly.  Equation \eqref{eq:p071-cross-collar-exact-collar} also shows
that $G_n$ equals the ambient affine function
$c+C_n+d_n+A(x)+B(y)$ on the fixed cross-collar.  This proves
\eqref{eq:p071-cross-collar-monotone-approximation} and the stated
first-stage membership.
\end{proof}
